% ─────────────────────────────────────────────────────────────────────────────
\chapter{Conclusions, Limitations and Future Work}
\label{chap:conclusions}
% ─────────────────────────────────────────────────────────────────────────────

% ─────────────────────────────────────────────────────────────────────────────
\section{Summary of the Research}
\label{sec:conclusions-summary}
% ─────────────────────────────────────────────────────────────────────────────

Il presente lavoro nasce dalla necessità di valutare in modo sistematico
l'affidabilità operativa di sistemi \gls{ai}-based impiegati nel triage forense
di immagini. In molti scenari applicativi, tali sistemi vengono ancora valutati
prevalentemente su dati puliti, bilanciati e semanticamente coerenti. Nei workflow
di \gls{df}, tuttavia, le immagini possono presentare caratteristiche molto più
eterogenee: possono essere ambigue, fuori distribuzione, ricompresse,
ridimensionate, alterate o deliberatamente manipolate. In questo contesto, la
sola prestazione nominale non è sufficiente per stabilire se un sistema possa
essere utilizzato in modo affidabile come supporto alla revisione investigativa.

La tesi ha affrontato questo problema adottando una prospettiva centrata sulla
robustezza operativa. Gli attacchi adversariali e le trasformazioni anti-forensi
non sono stati trattati come fine autonomo di \gls{advml}, ma come stressor
sperimentali utili a osservare il comportamento di modelli e tool in condizioni
non ideali. L'obiettivo principale è stato quindi verificare se l'automazione
\gls{ai}-based mantenga un comportamento coerente quando viene esposta a input
clean, campioni \gls{ood}, trasformazioni anti-forensi e perturbazioni
adversariali. Tale analisi ha incluso anche l'osservazione e la normalizzazione
degli output prodotti da strumenti commerciali black-box, al fine di confrontare
il comportamento dei modelli proxy con quello di sistemi utilizzati in scenari
forensi realistici.

Il percorso sperimentale ha compreso la costruzione di un dataset frozen da
1\,500 immagini, la valutazione dei modelli proxy EfficientNet-B0, ResNet18 e
\gls{clip} in condizioni clean, \gls{ood}, anti-forensi e adversariali, e la
produzione del forensic evaluation bundle da 11\,500 file per la valutazione
black-box di \gls{magnetaxiomtool}/\gls{magnetai}, \gls{excirefoto2025},
\gls{cellebriteinseyets} e \gls{griffeye}/\gls{t3kcore}.

Il protocollo sperimentale adottato è denominato \gls{fairlab}. Esso non
rappresenta un nuovo tool software né una nuova architettura generale, ma una
procedura riproducibile e auditabile per la valutazione della robustezza
operativa di sistemi \gls{ai}-based in scenari di \gls{df}. La sua struttura
consente di collegare dataset, modelli proxy, perturbazioni, trasformazioni,
bundle cieco, manifest, hash, output commerciali normalizzati, metriche
quantitative e analisi qualitativa \gls{xai} in un'unica pipeline tracciabile.
In questo modo, la coerenza tra costruzione del dataset, stress testing e
valutazione black-box non dipende da passaggi impliciti, ma da artefatti
verificabili e ricostruibili.

Il risultato complessivo della tesi è che la validazione di strumenti
\gls{ai}-based in ambito forense non può essere ridotta alla verifica della
prestazione nominale su dati puliti. Essa deve essere trattata come una procedura
sperimentale specifica per task, versione del tool, configurazione operativa e
condizioni di input. In questa prospettiva, la robustezza non è una proprietà
astratta del modello, ma una proprietà osservabile del sistema nel workflow
investigativo in cui viene impiegato.


% ─────────────────────────────────────────────────────────────────────────────
\section{Main Contributions}
\label{sec:conclusions-contributions}
% ─────────────────────────────────────────────────────────────────────────────

Il contributo della tesi si articola lungo quattro dimensioni complementari:
la definizione di un protocollo metodologico, la costruzione di una base
sperimentale controllata, la predisposizione di un benchmark tracciabile e
l'interpretazione dei risultati in chiave forense. Questi elementi non sono
indipendenti tra loro, ma concorrono a definire una pipeline coerente per
valutare la robustezza operativa di sistemi \gls{ai}-based impiegati nel triage
di immagini.

Sul piano metodologico, la tesi introduce \gls{fairlab} come protocollo
sperimentale integrato per la valutazione della robustezza operativa di sistemi
\gls{ai}-based in \gls{df}. Il valore del protocollo non risiede nella proposta
di un nuovo classificatore, ma nella formalizzazione di una procedura auditabile
che consente di testare sistemi esistenti in condizioni non ideali. In questa
prospettiva, dataset, modelli proxy, perturbazioni adversariali, trasformazioni
anti-forensi, metriche di valutazione, controlli di tracciabilità, output
commerciali normalizzati e analisi qualitativa vengono collegati in modo
esplicito all'interno di un unico workflow sperimentale.

La dimensione sperimentale del contributo riguarda la costruzione di una base
dati congelata e riproducibile. Il dataset finale comprende 1\,500 immagini
bilanciate nelle classi \texttt{weapon}, \texttt{non\_weapon} e \texttt{ood},
mentre il subset binario ufficiale comprende 1\,000 immagini destinate alla
generazione e alla valutazione delle perturbazioni. A partire da questa base sono
state prodotte cinque famiglie di perturbazioni adversariali, ossia \gls{fgsm},
\gls{superdeepfool}, \gls{sigmazero}, \gls{onepixel} e \gls{colorshift}, e
cinque trasformazioni anti-forensi model-agnostic, ossia
\gls{jpegrecompression}, \gls{resampleresize}, \gls{gaussianblur},
\gls{histogrammodification} e \gls{contraststretching}. Tale organizzazione ha
consentito di osservare il comportamento dei sistemi analizzati in condizioni
clean, fuori distribuzione, degradate e deliberatamente manipolate.

A questa base sperimentale si affianca la costruzione del forensic evaluation
bundle, composto da 11\,500 file organizzati in condizioni clean, \gls{ood},
adversariali e anti-forensi. Il bundle è stato predisposto per la valutazione
black-box mediante naming cieco, manifest separati, identificativi neutri e hash
\gls{sha256}/\gls{md5}. Questa struttura consente di collegare ogni output finale
al relativo artefatto sperimentale, mantenendo separati contenuto, condizione di
test, provenienza, identificativo e risultato metrico. In tal modo, la valutazione
dei tool commerciali non dipende da passaggi impliciti o da associazioni manuali
non documentate, ma da artefatti verificabili e ricostruibili.

La componente valutativa e interpretativa completa il contributo della tesi
attraverso il confronto tra tre modelli proxy, EfficientNet-B0, ResNet18 e
\gls{clip}, e i sistemi commerciali inclusi nella valutazione finale:
\gls{magnetaxiomtool}/\gls{magnetai}, \gls{excirefoto2025},
\gls{cellebriteinseyets} e \gls{griffeye}/\gls{t3kcore}. \gls{excirefoto2025} è
stato considerato come software general-purpose con funzionalità \gls{ai}-based
di semantic retrieval, non come prodotto esclusivamente forense; allo stesso
modo, \gls{cellebriteinseyets} è stato trattato come commercial black-box
\gls{ai}-assisted media analysis tool, mentre \gls{griffeye}/\gls{t3kcore} è
stato valutato attraverso i bookmark semantici automatici generati da
\gls{t3kcore}. Gli output osservabili dei tool sono stati esportati,
normalizzati e ricodificati operativamente rispetto al task
\texttt{weapon}/\texttt{non\_weapon}, mantenendo separata la valutazione dei
campioni \gls{ood}. L'analisi qualitativa tramite \gls{ig} su casi
rappresentativi di EfficientNet-B0 integra questa lettura, offrendo un supporto
diagnostico per interpretare alcuni failure mode senza sostituire le metriche
quantitative.


% ─────────────────────────────────────────────────────────────────────────────
\section{Main Experimental Findings}
\label{sec:conclusions-findings}
% ─────────────────────────────────────────────────────────────────────────────

I risultati sperimentali confermano che la robustezza operativa non può essere
valutata esclusivamente a partire dalla baseline clean. La prestazione nominale
su dati puliti rappresenta un punto di partenza necessario, ma non sufficiente
per stimare l'affidabilità di un sistema \gls{ai}-based in un workflow forense.
Nel contesto analizzato, infatti, modelli e tool con prestazioni aggregate elevate
possono manifestare comportamenti critici quando vengono esposti a input fuori
distribuzione, trasformazioni anti-forensi o perturbazioni adversariali.

Questa evidenza è particolarmente rilevante perché il triage forense non opera
su dati ideali. Le immagini acquisite durante un'indagine possono essere
ricompresse, ridimensionate, parzialmente degradate, semanticamente ambigue o
deliberatamente manipolate. Per questa ragione, la valutazione sperimentale ha
considerato più condizioni di input, osservando non solo la capacità dei sistemi
di classificare correttamente immagini clean, ma anche la loro stabilità rispetto
a condizioni realistiche o avversariali. La
Tabella~\ref{tab:conclusions-main-findings} sintetizza le principali evidenze
emerse, collegando ciascuna area sperimentale alla relativa implicazione
operativa.

\begin{table}[H]
\centering
\caption{Sintesi interpretativa delle principali evidenze sperimentali.}
\label{tab:conclusions-main-findings}
\footnotesize
\renewcommand{\arraystretch}{1.2}
\begin{tabularx}{\textwidth}{@{}p{0.18\textwidth} X X@{}}
\toprule
\textbf{Area} & \textbf{Evidenza principale} & \textbf{Rilevanza operativa} \\
\midrule
Clean baseline
& EfficientNet-B0 mostra la prestazione clean più elevata tra i proxy
  (0.978); tra i tool commerciali, \gls{griffeye}/\gls{t3kcore} ottiene la
  clean accuracy più elevata nel mapping conservativo adottato, mentre
  \gls{cellebriteinseyets} presenta un profilo bilanciato con recall
  \texttt{weapon} elevata.
& La prestazione nominale è necessaria per descrivere il comportamento di base
  del sistema, ma non è sufficiente per valutarne l'affidabilità in condizioni
  operative. \\

OOD
& \gls{clip} presenta il tasso più alto di predizioni \texttt{weapon} sui
  campioni \gls{ood} (0.8356); i modelli discriminativi mostrano invece livelli
  elevati di confidenza anche in assenza di contenuto pertinente al task.
& Gli input fuori distribuzione possono essere assorbiti nella classe
  \texttt{weapon}, generando rumore operativo, prioritizzazioni fuorvianti o
  carico aggiuntivo per l'analista. \\

Anti-forensic
& Le trasformazioni anti-forensi producono drop generalmente contenuti, ma
  introducono errori operativamente rilevanti, soprattutto falsi negativi su
  \texttt{weapon}.
& Manipolazioni realistiche come blur, ricompressione o alterazioni tonali
  possono incidere sul triage automatico anche quando il degrado metrico
  aggregato appare limitato. \\

Adversarial
& EfficientNet-B0 mostra la degradazione più marcata sotto \gls{sigmazero}
  (accuracy 0.015); nei tool commerciali black-box gli effetti delle
  perturbazioni risultano eterogenei e dipendono da sistema, configurazione e
  mapping operativo.
& Un'elevata accuracy clean non garantisce robustezza rispetto a input
  deliberatamente manipolati, soprattutto quando il sistema è esposto a
  perturbazioni progettate per alterarne la decisione. \\

Black-box tools
& I tool commerciali mostrano profili operativi differenziati:
  \gls{griffeye}/\gls{t3kcore} risulta più selettivo nel mapping conservativo
  adottato, \gls{cellebriteinseyets} mantiene un profilo bilanciato,
  \gls{magnetai} e \gls{excirefoto2025} evidenziano failure mode specifici.
& I tool commerciali devono essere interpretati attraverso output osservabili,
  soglie operative, mapping e normalizzazione black-box, evitando confronti
  diretti non controllati con i modelli proxy. \\

XAI
& Le mappe \gls{ig} rendono ispezionabili casi rappresentativi, ma non
  costituiscono prova autonoma di robustezza.
& L'\gls{xai} supporta la revisione critica e la diagnosi qualitativa dei
  failure mode, senza sostituire metriche quantitative o validazione esperta. \\
\bottomrule
\end{tabularx}
\end{table}

La lettura congiunta di questi risultati mostra una dissociazione significativa
tra accuracy nominale e affidabilità forense. In condizioni clean, alcuni sistemi
raggiungono prestazioni elevate e appaiono idonei al task considerato; tuttavia,
lo stesso comportamento non si conserva necessariamente quando cambiano la
distribuzione degli input, la qualità visiva dell'immagine o la natura della
manipolazione. Questo aspetto è centrale in un contesto investigativo, perché
l'errore non ha sempre lo stesso peso operativo: un falso negativo su
\texttt{weapon} può ridurre la probabilità che un contenuto rilevante venga
sottoposto a revisione umana, mentre un falso positivo o un flag improprio su
campioni \gls{ood} può aumentare il rumore operativo e alterare la
prioritizzazione del materiale da analizzare.

I risultati sui campioni \gls{ood} rafforzano questa interpretazione. Il fatto
che alcuni sistemi assegnino etichette \texttt{weapon} o livelli elevati di
confidenza anche a immagini non pertinenti al task indica che l'automazione può
produrre decisioni apparentemente sicure in assenza di contenuto
investigativamente rilevante. Questo comportamento non è semplicemente un errore
di classificazione nel senso tradizionale, ma un rischio operativo legato alla
gestione dell'incertezza e alla capacità del sistema di riconoscere i limiti del
proprio dominio applicativo.

Le trasformazioni anti-forensi mostrano un effetto meno estremo rispetto agli
attacchi adversariali più aggressivi, ma risultano comunque rilevanti sul piano
forense. Ricompressione \gls{jpeg}, blur, ridimensionamento e modifiche tonali sono
operazioni plausibili in scenari reali e possono essere introdotte sia da normali
processi di condivisione sia da condotte deliberatamente orientate a ridurre la
leggibilità automatica del contenuto. Anche quando il calo aggregato di accuracy
rimane contenuto, la comparsa di falsi negativi su immagini \texttt{weapon}
dimostra che la robustezza deve essere interpretata rispetto agli effetti sul
workflow investigativo, non soltanto rispetto alla variazione media delle
metriche.

Gli attacchi adversariali evidenziano invece il limite più netto della
prestazione nominale. Il caso di EfficientNet-B0 sotto \gls{sigmazero} mostra
che un modello molto accurato in condizioni clean può collassare in presenza di
perturbazioni costruite per alterarne la decisione. Nei tool commerciali
black-box, l'interpretazione è necessariamente più cauta, poiché non sono noti
architettura, soglie interne, strategie di pre-processing o criteri di retrieval.
Proprio per questo motivo, il confronto deve basarsi sugli output osservabili e
sulla loro ricodifica operativa, evitando di attribuire ai tool proprietà interne
non verificabili.

Nel complesso, la principale evidenza sperimentale della tesi è che la
validazione su dati puliti non equivale alla validazione in condizioni operative.
Un sistema può risultare accurato in baseline e, allo stesso tempo, produrre
falsi negativi su immagini \texttt{weapon}, flag impropri su campioni \gls{ood}
o errori dipendenti dalla soglia di retrieval. La robustezza operativa deve
quindi essere valutata come proprietà contestuale del sistema nel workflow
forense, considerando congiuntamente prestazioni clean, comportamento
\gls{ood}, sensibilità alle manipolazioni, output black-box e necessità di
revisione umana.



% ─────────────────────────────────────────────────────────────────────────────
\section{Operational and Forensic Implications}
\label{sec:conclusions-operational-implications}
% ─────────────────────────────────────────────────────────────────────────────

Le evidenze sperimentali hanno implicazioni dirette per l'uso responsabile dei
sistemi \gls{ai}-based nei workflow forensi. Il punto centrale riguarda il ruolo
che tali sistemi possono assumere nel triage investigativo. La classificazione
automatica può supportare la prioritizzazione delle evidenze e ridurre il carico
iniziale dell'analista, ma il suo output deve essere interpretato come
un'indicazione investigativa da verificare, contestualizzare e documentare, non
come una determinazione probatoria autonoma. La sua utilità dipende dalla
possibilità di ricostruire il contesto di generazione del risultato, la
configurazione del tool, il dataset di riferimento, le soglie operative adottate
e il successivo controllo umano. Questa esigenza è particolarmente rilevante per
i tool black-box, nei quali l'osservabilità del processo decisionale è limitata
agli output esportati e alla loro normalizzazione sperimentale.

In tale prospettiva, i falsi negativi assumono un rilievo operativo specifico.
In un workflow di triage forense, un falso negativo sulla classe
\texttt{weapon} può essere più critico di un falso positivo, poiché può ridurre
la probabilità che un contenuto investigativamente rilevante venga sottoposto a
revisione umana. Il rischio non riguarda soltanto i modelli proxy, ma anche i
tool commerciali, e può aumentare in presenza di perturbazioni o trasformazioni
capaci di alterare la leggibilità automatica dell'immagine. Il caso più severo
osservato nel protocollo è rappresentato da EfficientNet-B0 sotto
\gls{sigmazero}, con 489 transizioni
\texttt{weapon}\(\rightarrow\)\texttt{non\_weapon}. Questo risultato mostra in
modo evidente che un sistema accurato in condizioni clean può diventare
operativamente fragile quando viene esposto a input deliberatamente manipolati.

Un ulteriore elemento di cautela riguarda gli input \gls{ood}. Nel presente
lavoro, tali campioni non sono stati trattati come una normale terza classe
supervisionata, ma come un rischio operativo autonomo. Le immagini fuori
distribuzione possono infatti essere assorbite impropriamente nella classe
\texttt{weapon}, generare flag non pertinenti, aumentare il carico di revisione
dell'analista o produrre prioritizzazioni fuorvianti. In ambito forense, questo
comportamento è particolarmente delicato, perché un output apparentemente sicuro
su contenuti non appartenenti al dominio del task può indurre una falsa
percezione di affidabilità del sistema.

La configurazione operativa dei sistemi di semantic retrieval rappresenta un
altro fattore che incide direttamente sull'interpretazione dei risultati. Nel
caso di \gls{excirefoto2025}, il parametro \texttt{Distance limit} modifica il
trade-off tra falsi negativi e falsi positivi: la configurazione \texttt{D20}
risulta più restrittiva, \texttt{D80} più permissiva e \texttt{D50} costituisce
un compromesso intermedio nel contesto sperimentale. Ne deriva che la
configurazione del tool non può essere considerata un dettaglio tecnico
secondario, ma deve essere documentata come parte integrante del workflow
forense. A parità di contenuto analizzato, soglie diverse possono produrre
risultati differenti e incidere sulla successiva attività di revisione.

Queste considerazioni rafforzano l'importanza della tracciabilità e
dell'auditabilità dell'intero processo. La pipeline \gls{fairlab} mantiene
separati dataset, manifest, hash, bundle cieco, predizioni, metriche e output dei
tool, consentendo di ricostruire il percorso sperimentale e di distinguere tra
errori di classificazione, mancate rilevazioni, output non interpretabili e
limiti del mapping operativo. Tale separazione non ha soltanto valore tecnico,
ma costituisce una condizione metodologica essenziale per rendere l'uso forense
dell'\gls{ai} verificabile e difendibile. In assenza di artefatti tracciabili, il
risultato prodotto da un sistema automatico rischia di rimanere un'indicazione
opaca, difficilmente controllabile e scarsamente contestualizzabile.

Nel complesso, i risultati ottenuti confermano la necessità di mantenere un
paradigma \textit{human-in-the-loop} nell'impiego di strumenti \gls{ai}-based in
ambito forense. La classificazione automatica dovrebbe essere intesa come un
supporto al triage investigativo, da integrare con revisione umana,
documentazione delle operazioni e contestualizzazione del risultato. Tale
impostazione è coerente con i principi di supervisione umana, affidabilità e
accountability richiamati nelle linee guida europee per una \gls{ai} affidabile
\cite{EUAI2020}. Ciò vale in modo particolare per contenuti ad alta rilevanza
investigativa, input fuori distribuzione, immagini manipolate e output prodotti
da tool black-box, nei quali l'affidabilità operativa non può essere presunta,
ma deve essere valutata, documentata e sottoposta a controllo critico.



% ─────────────────────────────────────────────────────────────────────────────
\section{Limitations}
\label{sec:conclusions-limitations}
% ─────────────────────────────────────────────────────────────────────────────

Il presente lavoro presenta alcuni limiti che devono essere dichiarati
esplicitamente, affinché le conclusioni siano interpretate in modo corretto.
Tali limiti non riducono la validità del protocollo sperimentale, ma ne
delimitano il perimetro applicativo e chiariscono le condizioni entro cui i
risultati possono essere considerati rappresentativi.

Il principale vincolo interpretativo riguarda il dataset e il task sperimentale.
Il corpus costruito è finalizzato a un problema investigativo specifico, ossia
la classificazione binaria tra \texttt{weapon} e \texttt{non\_weapon}, con
valutazione separata dei campioni \gls{ood}. La dimensione complessiva, pari a
1\,500 immagini, consente una valutazione controllata e tracciabile, ma non
copre l'intera variabilità dei contenuti visivi che possono emergere in indagini
reali. Per questa ragione, i risultati non devono essere generalizzati
automaticamente a task multi-classe, domini visivi differenti, contenuti forensi
di altra natura o scenari nei quali la rilevanza investigativa dipenda da
elementi contestuali non rappresentati nel dataset.

La costruzione del corpus introduce inoltre una componente
\textit{human-in-the-loop}. Tale scelta aumenta la qualità semantica del dataset
e riduce il rischio di includere campioni ambigui o non pertinenti, ma comporta
anche una componente di giudizio umano nella selezione finale. Questa componente
non viene nascosta né trattata come automatica: è documentata mediante manifest,
criteri di inclusione, criteri di esclusione e bilanciamento. Di conseguenza, la
riproducibilità riguarda la procedura sperimentale congelata e gli artefatti
finali, mentre la fase di costruzione del dataset deve essere interpretata come
un processo controllato e auditabile, non come una procedura interamente
automatica.

A questo perimetro sperimentale si collega la copertura necessariamente
selettiva delle perturbazioni e delle trasformazioni considerate. Gli attacchi
adversariali e le trasformazioni anti-forensi selezionate non coprono l'intero
spazio delle manipolazioni possibili. La scelta è stata guidata dalla
rappresentatività operativa, dalla sostenibilità computazionale e dalla coerenza
con la finalità forense della tesi, non dall'esaustività. Gli attacchi
dipendenti dal modello sono generati rispetto a EfficientNet-B0 e devono quindi
essere interpretati alla luce dello scenario white-box considerato. Le
trasformazioni anti-forensi, invece, sono model-agnostic e rappresentano
degradazioni realistiche dell'immagine, ma non esauriscono tutte le possibili
tecniche di occultamento, manipolazione o contraffazione applicabili a evidenze
visive.

Un ulteriore elemento di cautela deriva dalla comparabilità tra modelli proxy e
strumenti commerciali black-box. I proxy model consentono una valutazione
controllata, con accesso a predizioni, confidenza e condizioni sperimentali
pienamente tracciabili. I software commerciali, invece, sono osservabili
soltanto attraverso gli output esportabili e non permettono di accedere ad
architettura, dati di addestramento, soglie decisionali, logica di
pre-processing o meccanismi interni di classificazione. Per questo motivo, il
confronto tra proxy e tool commerciali deve essere inteso come confronto
operativo di triage, non come valutazione di classificatori omogenei su label
identiche.

Anche la ricodifica delle categorie commerciali verso le classi
\texttt{weapon} e \texttt{non\_weapon} introduce una componente di incertezza
interpretativa. Il mapping adottato è controllato e documentato, ma non può
eliminare completamente i limiti derivanti dall'opacità dei sistemi black-box.
Questo aspetto è particolarmente rilevante per \gls{griffeye}/\gls{t3kcore},
dove il basso \gls{fpr} osservato riflette anche la scelta conservativa del
bookmark positivo primario e non può essere interpretato come prova autonoma di
superiorità intrinseca del tool.

Sul piano dell'interpretazione forense, occorre inoltre considerare che gli
output automatici non sono pienamente separabili dal contesto tecnico in cui
vengono prodotti. La verifica sui metadata sensibili ha mostrato che una parte
degli output commerciali può essere influenzata da informazioni incorporate nei
file. La costruzione del bundle cieco riduce questo rischio, ma non dimostra
l'assenza di qualunque forma di dipendenza residua da metadati, nomi interni o
artefatti non visivi. Questa cautela è particolarmente importante nei tool
black-box, nei quali non è possibile verificare direttamente quali segnali siano
stati utilizzati dal sistema per produrre una determinata classificazione o un
determinato risultato di retrieval.

Anche l'analisi \gls{xai} deve essere interpretata entro un perimetro
specifico. Le mappe basate su \gls{ig} sono state utilizzate come studio
qualitativo su casi rappresentativi di EfficientNet-B0 e supportano
l'interpretazione diagnostica dei failure case. Esse, tuttavia, non
costituiscono una metrica quantitativa di robustezza né una prova della
correttezza del modello. Il loro valore risiede nella possibilità di rendere più
ispezionabili alcuni comportamenti locali del classificatore, non nella
sostituzione delle metriche quantitative o della validazione esperta.

Infine, il presente lavoro non include una valutazione dell'ammissibilità
giuridica degli output prodotti dai sistemi \gls{ai}-based analizzati. Le
metriche adottate, come accuracy, weapon recall, \gls{fnr}, matching rate e
weapon flag rate, sono metriche operative di triage e non criteri di validazione
probatoria. L'output di un classificatore o di un tool di semantic retrieval non
costituisce, di per sé, una determinazione probatoria autonoma: la sua rilevanza
deve essere valutata nel contesto della catena di custodia digitale, della
documentazione della versione e della configurazione del tool, della validazione
indipendente e del quadro normativo applicabile.


% ─────────────────────────────────────────────────────────────────────────────
\section{Future Work}
\label{sec:conclusions-future-work}
% ─────────────────────────────────────────────────────────────────────────────

Le direzioni di sviluppo futuro derivano direttamente dai limiti dichiarati e
dalle evidenze emerse nel corso del lavoro. Il protocollo proposto offre una
base riproducibile per la valutazione della robustezza operativa, ma il suo
valore può essere rafforzato estendendone il perimetro sperimentale, migliorando
la gestione dell'incertezza e consolidando ulteriormente gli aspetti di
verificabilità e auditabilità.

Un'estensione naturale del lavoro consiste nell'applicare il protocollo a
ulteriori versioni, configurazioni e categorie operative dei tool commerciali
analizzati. Poiché \gls{magnetaxiomtool}/\gls{magnetai}, \gls{excirefoto2025},
\gls{cellebriteinseyets} e \gls{griffeye}/\gls{t3kcore} sono sistemi black-box o
quasi black-box, i risultati ottenuti su una versione specifica non devono essere
assunti come validi per release future, impostazioni differenti o funzionalità
non esportate nel presente studio. Una validazione più ampia dovrebbe quindi
ripetere il protocollo al variare delle versioni software, delle configurazioni
operative e delle categorie semantiche disponibili, così da distinguere tra
proprietà stabili del sistema e comportamenti dipendenti dalla specifica
configurazione esaminata.

L'ampliamento della validazione dovrebbe riguardare anche il dataset e i task
sperimentali. Il lavoro potrebbe essere esteso verso scenari multi-classe,
domini visivi differenti e immagini acquisite in condizioni operative più
eterogenee, includendo export da dispositivi reali, contenuti provenienti da
piattaforme social o app di messaggistica e immagini degradate da pipeline di
acquisizione forense. In questa prospettiva, sarebbe utile integrare procedure
di revisione multi-annotatore o misure di agreement sulla selezione semantica
dei campioni, in modo da quantificare la componente soggettiva introdotta dalla
validazione \textit{human-in-the-loop} e rendere ancora più trasparente il
processo di costruzione del corpus.

I risultati ottenuti sui campioni \gls{ood} suggeriscono inoltre la necessità di
approfondire la gestione dell'incertezza. Nel protocollo sperimentale, diversi
sistemi tendono ad assorbire contenuti anomali in una delle classi operative,
talvolta con livelli di confidenza non trascurabili. Una prosecuzione del lavoro
potrebbe quindi integrare meccanismi di uncertainty estimation, rejection option,
calibrazione della confidenza e \gls{ood} detection, con l'obiettivo di ridurre
il rischio di classificazioni forzate su input non pertinenti al task. In un
workflow forense, infatti, la capacità di riconoscere quando un contenuto non
appartiene al dominio decisionale previsto è rilevante quanto la capacità di
classificare correttamente i campioni appartenenti alle classi note.

Sempre in relazione alla robustezza, ulteriori sviluppi potrebbero considerare
strategie di attacco e difesa non incluse nel presente protocollo, come attacchi
transfer-based, adversarial patches, trasformazioni fisiche o ibride, forme di
adversarial training e tecniche di pre-processing robusto. Tali estensioni
dovrebbero tuttavia rimanere subordinate alla finalità forense della valutazione.
L'obiettivo non dovrebbe essere trasformare il lavoro in un benchmark puramente
algoritmico di \gls{advml}, ma comprendere in che misura specifiche condizioni
avversariali o anti-forensi incidano sull'affidabilità operativa dei sistemi
utilizzati a supporto del triage investigativo.

Un ulteriore asse di sviluppo riguarda la verificabilità del protocollo e la
disponibilità degli artefatti necessari alla riproduzione dell'esperimento.
Anche quando le immagini non possono essere redistribuite integralmente per
ragioni di licenza o sensibilità del dato, script, manifest, configurazioni,
hash, metriche aggregate e procedure di normalizzazione possono essere resi
disponibili per consentire una verifica indipendente della procedura
sperimentale. Questa esigenza è particolarmente rilevante per i tool commerciali
black-box, per i quali la ripetibilità della procedura e la chiarezza dei criteri
di mapping diventano elementi centrali della validazione.

In prospettiva, anche l'integrazione di metodi \gls{xai} nei software forensi
potrebbe contribuire a migliorare la trasparenza degli output automatici. Tale
integrazione dovrebbe però essere interpretata con cautela: le spiegazioni
visuali o testuali possono supportare la diagnosi dei failure mode e la revisione
critica da parte dell'analista, ma non devono essere trattate come prova autonoma
della correttezza della classificazione. Il loro valore risiede nella capacità di
rendere più ispezionabile il comportamento del sistema, non nella sostituzione
delle metriche quantitative, della validazione indipendente o del controllo
esperto.

Più in generale, il lavoro suggerisce la necessità di protocolli condivisi per
la validazione dell'\gls{ai} in \gls{df}. Tali protocolli dovrebbero includere
metriche orientate al contesto forense, versioning dei tool, documentazione delle
configurazioni, analisi degli errori, benchmark controllati e procedure di audit
indipendenti. Solo attraverso una validazione sistematica, tracciabile e
ripetibile è possibile integrare strumenti \gls{ai}-based nei workflow forensi in
modo metodologicamente difendibile, mantenendo il ruolo dell'automazione come
supporto al triage e non come sostituto della valutazione investigativa ed
esperta.


% ─────────────────────────────────────────────────────────────────────────────
\section{Closing Remarks}
\label{sec:conclusions-closing}
% ─────────────────────────────────────────────────────────────────────────────

La presente tesi ha mostrato che la robustezza operativa dei sistemi
\gls{ai}-based in \gls{df} non può essere ridotta a una misura di accuracy su
dati puliti. In presenza di input \gls{ood}, trasformazioni anti-forensi e
perturbazioni adversariali, sistemi con prestazioni nominali elevate possono
produrre errori rilevanti, anche ad alta confidenza, con implicazioni dirette per
il triage investigativo.

Il contributo metodologico del lavoro consiste nell'aver reso questo problema
sperimentalmente verificabile e tracciabile attraverso il protocollo
\gls{fairlab} e il forensic evaluation bundle. Il passaggio da una valutazione
astratta della robustezza a una valutazione orientata alle condizioni operative
del triage forense rappresenta il principale cambiamento di prospettiva proposto
dalla tesi. In tale prospettiva, la robustezza non viene interpretata come una
proprietà generale e statica del modello, ma come una proprietà osservabile del
sistema nel contesto operativo in cui viene impiegato.

Nel perimetro sperimentale considerato, i risultati mostrano che nessuno dei
sistemi valutati può essere considerato operativamente affidabile in modo
incondizionato. Ciascun sistema presenta
failure mode specifici, dipendenti dalla condizione di input, dalla
configurazione, dal task e dal tipo di output osservabile. Questa evidenza non
invalida l'uso dell'\gls{ai} nel workflow forense, ma ne definisce i prerequisiti
metodologici: validazione task-specific, documentazione della configurazione
operativa, gestione esplicita degli input anomali, tracciabilità degli artefatti
e controllo umano sul risultato.

Il protocollo \gls{fairlab} offre una struttura sperimentale per rendere tali
prerequisiti verificabili, auditabili e riproducibili. In questo senso, il lavoro
non propone l'automazione come sostituto dell'analista, ma come componente di
supporto che deve essere validata, contestualizzata e mantenuta all'interno di un
workflow forense controllato. La robustezza operativa diventa quindi una
condizione da dimostrare sperimentalmente, non una proprietà da presumere sulla
base delle sole prestazioni nominali.
